% !TeX root = ../../ONMIPA Alin.tex
% Sumber utama: Fuzhen Zhang, Linear Algebra: Challenging Problems for
% Students, edisi kedua, Bab 2; mulai hlm. cetak 21 / halaman PDF 40.
\chapter{Determinan}
\label{ch:determinan}

Pembahasan ini mengikuti bagian \emph{Definitions and Facts} pada
\citet[hlm.~21--23]{zhang2009}, yaitu halaman PDF 40--42, kemudian
mengembangkan teknik determinan dari soal-soal Bab 2 buku tersebut.
Operasi matriks, operasi baris elementer, dan pengertian invers dianggap
sudah dikenal. Semua matriks berentri dalam $\F=\R$ atau $\C$, kecuali
dinyatakan lain. Notasi $|A|$ berarti $\det A$; untuk suatu skalar
kompleks $z$, notasi $|z|$ tetap berarti modulus.

\section{Definisi dan Ekspansi}
\label{det:ekspansi}

\begin{defi}[Minor, kofaktor, dan determinan]
      Untuk $A=(a_{ij})\in\Mat_n(\F)$, tuliskan $A_{ij}$ untuk submatriks
      yang diperoleh dengan menghapus baris ke-$i$ dan kolom ke-$j$.
      Determinan didefinisikan secara rekursif oleh
      \[
            \det(a_{11})=a_{11},\qquad
            \det A=\sum_{j=1}^n(-1)^{1+j}a_{1j}\det A_{1j}\quad(n\ge2).
      \]
      Bilangan $M_{ij}=\det A_{ij}$ disebut \emph{minor} dari entri $a_{ij}$,
      sedangkan $C_{ij}=(-1)^{i+j}M_{ij}$ disebut \emph{kofaktornya}.
      Untuk konsistensi pada ukuran satu, determinan matriks kosong ditetapkan $1$.
\end{defi}

Definisi di atas menggunakan baris pertama. Untuk membuktikan bahwa
baris atau kolom lain memberikan nilai yang sama, diperlukan hubungan
dengan permutasi.

\begin{defi}[Tanda permutasi]
      Tuliskan $S_n$ untuk himpunan seluruh permutasi $\{1,\ldots,n\}$.
      Untuk $\sigma\in S_n$, definisikan banyaknya inversi
      \[
            \nu(\sigma)=\#\{(i,j):i<j,\ \sigma(i)>\sigma(j)\},
            \qquad \operatorname{sgn}(\sigma)=(-1)^{\nu(\sigma)}.
      \]
      Permutasi disebut genap jika tandanya $1$, dan ganjil jika tandanya $-1$.
      Satu pertukaran dua posisi membalik tanda permutasi.
\end{defi}

\begin{thm}[Rumus Leibniz]
      \label{det:leibniz}
      Untuk $A\in\Mat_n(\F)$,
      \[
            \det A=\sum_{\sigma\in S_n}\operatorname{sgn}(\sigma)
            \prod_{i=1}^n a_{i,\sigma(i)}.
      \]
      Setiap suku mengambil tepat satu entri dari setiap baris dan setiap kolom.
\end{thm}
\begin{proof}
      Gunakan induksi pada $n$. Kelompokkan permutasi menurut nilai
      $\sigma(1)=j$. Setelah posisi pertama dan nilai $j$ dihapus,
      permutasi sisanya menjadi permutasi pada $n-1$ unsur.
      Pasangan inversi yang melibatkan posisi pertama berjumlah $j-1$;
      karena $(-1)^{j-1}=(-1)^{1+j}$, kontribusi kelompok ini adalah
      $(-1)^{1+j}a_{1j}\det A_{1j}$.
      Jumlah seluruh kelompok tepat menghasilkan definisi rekursif.
\end{proof}

\begin{thm}[Ekspansi Laplace]
      \label{det:laplace}
      Untuk setiap baris $i$ dan setiap kolom $j$,
      \[
            \det A=\sum_{k=1}^n a_{ik}C_{ik}
            =\sum_{k=1}^n a_{kj}C_{kj}.
      \]
\end{thm}
\begin{proof}
      Kelompokkan suku Leibniz menurut entri yang diambil dari baris $i$.
      Pada kelompok yang mengambil entri $(i,k)$, menghapus baris $i$ dan
      kolom $k$ menyisakan determinan $A_{ik}$ dengan faktor tanda
      $(-1)^{i+k}$. Ini memberikan ekspansi baris.
      Pengelompokan menurut kolom memberikan ekspansi kolom dengan alasan sama.
\end{proof}

Rumus Leibniz terutama berguna untuk membaca struktur: jika tidak ada
permutasi yang memilih seluruh entri tak nol, determinannya nol.
Ekspansi Laplace berguna ketika terdapat baris atau kolom dengan sedikit
entri tak nol. Mengembangkan seluruh $n!$ suku biasanya bukan cara
perhitungan yang efektif.

\section{Sifat Fundamental}
\label{det:sifat}

\begin{thm}[Multilinear dan berselang-seling]
      \label{det:multilinear}
      Determinan linear terhadap \emph{masing-masing} kolom ketika kolom
      lain dipertahankan. Menukar dua kolom mengubah tanda determinan;
      jika dua kolom sama, determinannya nol. Pernyataan yang sama berlaku
      untuk baris.
\end{thm}
\begin{proof}
      Dalam setiap suku Leibniz terdapat tepat satu faktor dari kolom yang
      sedang diubah; distributivitas memberi linearitas terhadap kolom itu.
      Menukar dua kolom memasangkan setiap permutasi dengan permutasi
      bertanda berlawanan. Jika kolom yang ditukar sama, matriks tidak
      berubah, sehingga $\det A=-\det A$ dan $\det A=0$.
      Argumen baris identik.
\end{proof}

Akibatnya, tiga operasi elementer mempunyai pengaruh berikut:
\[
      \begin{array}{c|c}
            \text{Operasi pada }A\text{ menghasilkan }B & \det B  \\ \hline
            R_i\leftrightarrow R_j                      & -\det A \\
            R_i\leftarrow cR_i                          & c\det A \\
            R_i\leftarrow R_i+cR_j,\quad i\ne j         & \det A
      \end{array}
\]
Pada operasi kedua, $c\ne0$ diperlukan jika operasi hendak dibalik;
identitas determinannya sendiri juga benar untuk $c=0$.
Aturan yang sama berlaku untuk kolom. Aturan penambahan baris mengikuti
multilinearitas: suku tambahan mempunyai dua baris yang sama.

\begin{lem}[Karakterisasi determinan]
      \label{det:unik}
      Jika fungsi $f(v_1,\ldots,v_n)$ multilinear dan berselang-seling
      terhadap $n$ vektor kolom dalam $\F^n$, maka
      \[
            f(v_1,\ldots,v_n)=f(e_1,\ldots,e_n)\det[v_1\ \cdots\ v_n].
      \]
\end{lem}
\begin{proof}
      Ekspansikan setiap $v_j$ sebagai kombinasi vektor koordinat $e_i$.
      Multilinearitas menghasilkan jumlah nilai $f(e_{i_1},\ldots,e_{i_n})$.
      Suku dengan indeks berulang nol, sedangkan suku tanpa indeks berulang
      merupakan permutasi dari $f(e_1,\ldots,e_n)$ dengan tanda permutasinya.
      Jumlah koefisien yang tersisa adalah rumus Leibniz.
\end{proof}

\begin{thm}[Determinan hasil kali]
      \label{det:produk}
      Untuk $A,B\in\Mat_n(\F)$, berlaku $\det(AB)=\det A\det B$.
\end{thm}
\begin{proof}
      Dengan $A$ tetap, definisikan
      $f(b_1,\ldots,b_n)=\det[Ab_1\ \cdots\ Ab_n]$.
      Fungsi ini multilinear dan berselang-seling, serta
      $f(e_1,\ldots,e_n)=\det A$.
      Terapkan \autoref{det:unik} pada kolom-kolom $B$.
\end{proof}

\begin{prop}[Sifat-sifat dasar Zhang]
      \label{det:dasar}
      Untuk $A,B\in\Mat_n(\F)$, skalar $c$, dan matriks invertibel $S$:
      \begin{enumerate}[label=(d\arabic*),nosep]
            \item $A$ singular jika dan hanya jika $\det A=0$.
            \item Baris-baris atau kolom-kolom $A$ bergantung linear jika dan
                  hanya jika $\det A=0$.
            \item Jika $A$ segitiga, $\det A=\prod_{i=1}^n a_{ii}$.
            \item $\det(A^T)=\det A$.
            \item $\det(cA)=c^n\det A$.
            \item $\det(AB)=\det A\det B$.
            \item $\det(S^{-1}AS)=\det A$.
      \end{enumerate}
      \sumber{\citet[hlm.~22, sifat (d1)--(d7)]{zhang2009}}
\end{prop}
\begin{proof}
      Pada matriks segitiga, hanya permutasi identitas yang dapat memberi
      suku tak nol, sehingga (d3) berlaku. Mengganti permutasi dengan
      inversnya dalam rumus Leibniz membuktikan (d4); multilinearitas
      memberi (d5). Sifat (d6) sudah dibuktikan.

      Untuk (d1)--(d2), reduksi baris ke bentuk eselon mempertahankan
      kenolannya determinan. Jika terdapat baris nol, determinannya nol
      dan kolom-kolomnya tidak bebas. Jika semua pivot ada, matriks eselon
      segitiga mempunyai diagonal tak nol, sehingga determinannya tak nol
      dan matriks dapat direduksi menjadi $I_n$; semua OBE dapat dibalik,
      maka $A$ invertibel. Transpose memberikan pernyataan untuk baris.

      Akhirnya, $\det(S^{-1})\det S=\det I=1$. Dengan (d6),
      \[
            \det(S^{-1}AS)=\det(S^{-1})\det A\det S=\det A.
      \]
\end{proof}

\begin{exam}[Linear pada kolom, bukan pada matriks]
      Misalkan $A=[a\ r_2\ r_3\ r_4]$, $B=[b\ r_2\ r_3\ r_4]$,
      $\det A=4$, dan $\det B=1$. Maka
      \[
            \det(A+B)=2^3\det[a+b\ r_2\ r_3\ r_4]=8(4+1)=40.
      \]
      Jika $C=[r_4\ r_3\ r_2\ a+b]$, urutan empat kolom dibalik melalui
      $\binom42=6$ pertukaran, sehingga $\det C=(-1)^6(4+1)=5$.
      Pada umumnya $\det(A+B)\ne\det A+\det B$.
      \sumber{\citet[Soal 2.13]{zhang2009}}
\end{exam}

\subsection{Tanda permutasi dan kenolan struktural}
Jika urutan $n$ kolom dibalik, faktor tandanya adalah
$(-1)^{n(n-1)/2}$. Karena itu, untuk matriks dengan entri nol
di sebelah kiri diagonal sekundernya,
\[
      a_{ij}=0\text{ jika }i+j<n+1
      \quad\Longrightarrow\quad
      \det A=(-1)^{n(n-1)/2}\prod_{i=1}^n a_{i,n+1-i}.
\]
Pembalikan kolom mengubah matriks tersebut menjadi segitiga bawah.
Argumen ini juga menerangkan tanda determinan matriks pembalik urutan
pada \citet[Soal 2.31]{zhang2009}.

\begin{exam}[Determinan nol tanpa ekspansi]
      Determinan matriks berikut nol untuk semua nilai entrinya:
      \[
            \begin{vmatrix}
                  a_1 & a_2 & a_3 & a_4 & a_5 \\b_1&b_2&b_3&b_4&b_5\\
                  c_1 & c_2 & 0   & 0   & 0   \\d_1&d_2&0&0&0\\e_1&e_2&0&0&0
            \end{vmatrix}=0.
      \]
      Dalam suku Leibniz yang tidak nol, tiga baris terakhir harus memilih
      kolom dari himpunan $\{1,2\}$. Tiga kolom yang berbeda tidak dapat
      dipilih dari dua kemungkinan. Jadi setiap suku Leibniz nol.
      \sumber{\citet[Soal 2.3]{zhang2009}}
\end{exam}

\section{Adjugat dan Perubahan Entri}
\label{det:adjugat}

\begin{defi}[Adjugat]
      Matriks adjugat $A$, ditulis $\operatorname{adj}(A)$, adalah transpose
      matriks kofaktor:
      \[
            (\operatorname{adj}(A))_{ij}=C_{ji}.
      \]
      Istilah \emph{adjoint} pada halaman 22 buku Zhang menunjuk ke matriks
      ini, bukan transpose konjugat $A^*=\overline A^{\,T}$.
\end{defi}

\begin{thm}[Identitas adjugat]
      \label{det:identitas-adj}
      Untuk setiap matriks persegi, termasuk yang singular,
      \[
            A\operatorname{adj}(A)=\operatorname{adj}(A)A=(\det A)I_n.
      \]
      Jika $\det A\ne0$, maka $A^{-1}=(\det A)^{-1}\operatorname{adj}(A)$.
\end{thm}
\begin{proof}
      Entri $(i,j)$ dari $A\operatorname{adj}(A)$ adalah
      $\sum_k a_{ik}C_{jk}$. Jika $i=j$, jumlah ini adalah ekspansi Laplace
      dan bernilai $\det A$. Jika $i\ne j$, jumlah tersebut adalah
      determinan matriks yang baris ke-$j$-nya diganti dengan baris ke-$i$.
      Kofaktor pada baris $j$ tidak berubah karena baris itu dihapus ketika
      minor dihitung. Matriks pengganti mempunyai dua baris sama, sehingga
      determinannya nol. Identitas dari sisi lain dibuktikan dengan kolom.
\end{proof}

\begin{thm}[Perubahan rank satu]
      \label{det:rankone}
      Untuk $u,v\in\F^n$ dan setiap $A\in\Mat_n(\F)$,
      \[
            \boxed{\det(A+uv^T)=\det A+v^T\operatorname{adj}(A)u.}
      \]
      Jika $A$ invertibel, rumus ini menjadi
      \[
            \det(A+uv^T)=\det A\,(1+v^TA^{-1}u).
      \]
\end{thm}
\begin{proof}
      Kolom ke-$j$ dari $A+uv^T$ adalah $a_j+v_ju$.
      Ekspansikan determinan secara multilinear. Semua suku yang memilih
      $u$ dari dua kolom atau lebih nol, karena dua kolom itu sebanding.
      Suku tanpa $u$ adalah $\det A$; jumlah suku dengan tepat satu $u$ ialah
      \[
            \sum_{j=1}^n v_j\sum_{i=1}^n u_iC_{ij}
            =v^T\operatorname{adj}(A)u.
      \]
      Untuk $A$ invertibel, gunakan
      $\operatorname{adj}(A)=(\det A)A^{-1}$.
\end{proof}

Khususnya, $\det(I_n+uv^T)=1+v^Tu$. Dengan $v^T$ diganti oleh
$-y^*$, diperoleh $\det(I_n-xy^*)=1-y^*x$, yaitu identitas dalam
\citet[Soal 2.41(a)]{zhang2009}. Rumus adjugat di atas tetap sah
ketika $A$ singular; rumus yang mengandung $A^{-1}$ tidak boleh
dipakai pada kasus tersebut.

\begin{prop}[Determinan sebagai polinom]
      \label{det:principal}
      Jika $A[S,S]$ menyatakan submatriks utama pada himpunan indeks
      $S\subseteq\{1,\ldots,n\}$, maka
      \[
            \det(I_n+tA)=\sum_{S\subseteq\{1,\ldots,n\}}
            t^{|S|}\det A[S,S].
      \]
      Suku untuk $S=\varnothing$ bernilai $1$. Khususnya, koefisien $t$
      adalah $\tr(A)$ dan koefisien $t^n$ adalah $\det A$.
\end{prop}
\begin{proof}
      Pada kolom ke-$j$, pilih $e_j$ atau $ta_j$ lalu gunakan multilinearitas.
      Jika $ta_j$ dipilih tepat untuk $j\in S$, kolom di luar $S$ memaksa
      permutasi Leibniz menetapkan setiap indeks di luar $S$.
      Jumlah suku yang tersisa adalah $t^{|S|}\det A[S,S]$.
\end{proof}

Untuk $n\ge2$, koefisien $t^2$ adalah
\[
      \sum_{i<j}(a_{ii}a_{jj}-a_{ij}a_{ji})
      =\frac12\bigl((\tr A)^2-\tr(A^2)\bigr).
\]
Dengan demikian, beberapa informasi determinan dapat diperoleh tanpa
menghitung keseluruhan polinom.

\begin{prop}[Rumus diferensiasi Jacobi]
      Jika setiap entri $A(t)$ dapat diturunkan terhadap parameter real $t$,
      maka
      \[
            \frac{d}{dt}\det A(t)
            =\sum_{i,j}C_{ij}(t)a'_{ij}(t)
            =\tr\bigl(\operatorname{adj}(A(t))A'(t)\bigr).
      \]
      Jika $A(t)$ invertibel, bentuknya menjadi
      $\det A(t)\,\tr(A(t)^{-1}A'(t))$.
      \sumber{pengembangan rumus diferensiasi per kolom pada
            \citet[Soal 2.11]{zhang2009}}
\end{prop}
\begin{proof}
      Turunkan rumus Leibniz dengan aturan turunan hasil kali.
      Kelompokkan suku menurut entri yang diturunkan; koefisien
      $a'_{ij}(t)$ adalah kofaktor $C_{ij}(t)$.
      Identitas trace mengikuti definisi adjugat.
\end{proof}
Sebagai akibat, $\left.\frac{d}{dt}\det(I_n+tA)\right|_{t=0}=\tr(A)$.
Rumus tanpa invers tetap berlaku pada nilai $t$ yang membuat $A(t)$ singular.

\section{Pola Entri dan Parameter}
\label{det:parameter}

\begin{thm}[Diagonal ditambah matriks rank satu]
      \label{det:diagrankone}
      Untuk $D=\diag(d_1,\ldots,d_n)$ dan $u,v\in\F^n$,
      \[
            \det(D+uv^T)=\prod_{i=1}^n d_i
            +\sum_{i=1}^n u_iv_i\prod_{j\ne i}d_j.
      \]
      Tidak disyaratkan semua $d_i$ tak nol.
\end{thm}
\begin{proof}
      Adjugat matriks diagonal adalah matriks diagonal dengan entri ke-$i$
      $\prod_{j\ne i}d_j$. Terapkan \autoref{det:rankone}.
\end{proof}

\begin{exam}[Entri di luar diagonal sama]
      Tentukan determinan matriks dengan entri diagonal $p_1,\ldots,p_n$
      dan seluruh entri di luar diagonal sama dengan $a$.
      \sumber{\citet[Soal 2.6(b)]{zhang2009}, ditulis dalam bentuk ekuivalen}
\end{exam}
\begin{solution}
      Dengan $\boldsymbol{1}=(1,\ldots,1)^T$, matriksnya adalah
      \[
            M=\diag(p_1-a,\ldots,p_n-a)+a\boldsymbol{1}\boldsymbol{1}^T.
      \]
      Maka
      \[
            \det M=\prod_{i=1}^n(p_i-a)
            +a\sum_{i=1}^n\prod_{j\ne i}(p_j-a).
      \]
      Jika $f(x)=\prod_i(p_i-x)$ dan $f_i(a)=\prod_{j\ne i}(p_j-a)$,
      hasil ini sama dengan $f(a)+a\sum_i f_i(a)$.
      Karena $f(a)=(p_n-a)f_n(a)$, bentuk yang dicantumkan Zhang ialah
      $a\sum_{i=1}^{n-1}f_i(a)+p_nf_n(a)$.

      Khususnya, jika semua entri diagonal adalah $\alpha$ dan semua entri
      di luar diagonal adalah $\beta$, maka
      \[
            \boxed{\det M=(\alpha-\beta)^{n-1}
                  \bigl(\alpha+(n-1)\beta\bigr).}
      \]
      Inilah hasil pada Soal 2.6(c). Untuk $n\ge2$, kedua kemungkinan
      singularnya adalah $\alpha=\beta$ atau $\alpha=-(n-1)\beta$.
\end{solution}

\begin{exam}[Dua konstanta di atas dan di bawah diagonal]
      \label{det:zhang26}
      Misalkan
      \[
            M=\begin{pmatrix}
                  p_1    & a      & a      & \cdots & a      \\
                  b      & p_2    & a      & \cdots & a      \\
                  b      & b      & p_3    & \cdots & a      \\
                  \vdots & \vdots & \vdots & \ddots & \vdots \\
                  b      & b      & b      & \cdots & p_n
            \end{pmatrix},\qquad f(x)=\prod_{i=1}^n(p_i-x).
      \]
      Untuk $a\ne b$, buktikan
      \[
            \det M=\frac{bf(a)-af(b)}{b-a}.
      \]
      \sumber{\citet[Soal 2.6(a)]{zhang2009}}
\end{exam}
\begin{solution}
      \ide{Kurangkan satu matriks rank satu agar sisanya segitiga.}
      Tuliskan $M=L+a\boldsymbol{1}\boldsymbol{1}^T$, dengan $L$ segitiga
      bawah, diagonal $p_i-a$, dan semua entri di bawah diagonal $b-a$.
      Mula-mula anggap $p_i\ne a$ untuk setiap $i$, sehingga $L$ invertibel.
      Jika $y=L^{-1}\boldsymbol{1}$ dan $s_i=y_1+\cdots+y_i$, $s_0=0$,
      persamaan baris ke-$i$ memberikan
      \[
            (p_i-a)y_i+(b-a)s_{i-1}=1.
      \]
      Definisikan $r_i=1-(b-a)s_i$. Karena
      \[
            r_i=r_{i-1}-(b-a)y_i
            =\frac{p_i-b}{p_i-a}r_{i-1},\qquad r_0=1,
      \]
      diperoleh $r_n=f(b)/f(a)$ dan $s_n=(1-r_n)/(b-a)$.
      Sekarang terapkan rumus perubahan rank satu:
      \begin{align*}
            \det M
             & =\det L\,(1+a\boldsymbol{1}^TL^{-1}\boldsymbol{1}) \\
             & =f(a)(1+as_n)
            =\frac{bf(a)-af(b)}{b-a}.
      \end{align*}
      Setelah dikalikan $b-a$, kedua ruas adalah polinom dalam
      $p_1,\ldots,p_n$. Identitas yang benar untuk semua $p_i\ne a$
      tetap benar ketika sebagian $p_i=a$, melalui identitas polinom
      atau kontinuitas. Jadi syarat sementara tersebut dapat dihapus.
      Kasus $a=b$ ditangani oleh contoh sebelumnya, bukan dengan membagi oleh nol.
\end{solution}

\section{Determinan Vandermonde}
\label{det:vandermonde}

\begin{thm}[Vandermonde]
      \label{det:rumus-vandermonde}
      Untuk $x_1,\ldots,x_n\in\F$,
      \[
            V_n(x_1,\ldots,x_n)=
            \begin{vmatrix}
                  1         & 1         & \cdots & 1         \\x_1&x_2&\cdots&x_n\\
                  x_1^2     & x_2^2     & \cdots & x_n^2     \\
                  \vdots    & \vdots    &        & \vdots    \\
                  x_1^{n-1} & x_2^{n-1} & \cdots & x_n^{n-1}
            \end{vmatrix}
            =\prod_{1\le i<j\le n}(x_j-x_i).
      \]
      \sumber{\citet[Soal 2.7]{zhang2009}}
\end{thm}
\begin{proof}
      Pandang $V_n$ sebagai polinom dalam $x_n$, dengan koefisien berupa polinom
      dalam $x_1,\ldots,x_{n-1}$. Derajatnya paling tinggi $n-1$.
      Jika $x_n=x_i$ untuk $i<n$, dua kolom sama, sehingga polinom tersebut
      memiliki faktor $x_n-x_i$. Koefisien $x_n^{n-1}$ adalah
      $V_{n-1}(x_1,\ldots,x_{n-1})$, hasil ekspansi pada entri $(n,n)$.
      Karena itu
      \[
            V_n=V_{n-1}\prod_{i=1}^{n-1}(x_n-x_i).
      \]
      Induksi dengan $V_1=1$ menyelesaikan bukti sebagai identitas polinom,
      termasuk saat beberapa $x_i$ berimpit.
\end{proof}

\begin{prop}[Basis polinom lain]
      Jika $p_i(x)$ berderajat tepat $i$ dengan koefisien utama $c_i$,
      untuk $i=0,\ldots,n-1$, maka
      \[
            \det\bigl(p_{i-1}(x_j)\bigr)_{i,j=1}^n
            =\left(\prod_{i=0}^{n-1}c_i\right)
            \prod_{1\le i<j\le n}(x_j-x_i).
      \]
\end{prop}
\begin{proof}
      Tuliskan setiap $p_i$ sebagai kombinasi $1,x,\ldots,x^i$.
      Matriks evaluasinya merupakan hasil kali matriks koefisien segitiga
      bawah dengan matriks Vandermonde. Diagonal matriks koefisien adalah
      $c_0,\ldots,c_{n-1}$; gunakan determinan hasil kali.
\end{proof}

\begin{exam}[Pangkat dan koefisien binomial]
      Untuk $V=(j^{i-1})_{i,j=1}^n$, rumus Vandermonde memberi
      \[
            \det V=\prod_{1\le i<j\le n}(j-i)
            =\prod_{j=2}^n(j-1)!=\prod_{k=1}^{n-1}k!.
      \]
      Sebagai pengembangan, jika $B_{ij}=\binom{j-1}{i-1}$, maka
      $\det B=1$. Memang, polinom
      \[
            p_k(x)=\frac{x(x-1)\cdots(x-k+1)}{k!},\qquad p_0(x)=1,
      \]
      mempunyai koefisien utama $1/k!$. Evaluasi pada $0,1,\ldots,n-1$
      menghasilkan faktor Vandermonde $\prod_{k=1}^{n-1}k!$, yang tepat
      membatalkan hasil kali koefisien utama.
      \sumber{bagian pertama \citet[Soal 2.7]{zhang2009}; bagian binomial
            merupakan pengembangan dari identitas yang sama}
\end{exam}

\section{Determinan melalui Rekurensi}
\label{det:rekurensi}

\begin{prop}[Matriks tridiagonal]
      Definisikan $D_n=\det T_n$, dengan
      \[
            T_n=\begin{pmatrix}
                  p      & q      & 0      & \cdots & 0      \\
                  r      & p      & q      & \ddots & \vdots \\
                  0      & r      & p      & \ddots & 0      \\
                  \vdots & \ddots & \ddots & \ddots & q      \\
                  0      & \cdots & 0      & r      & p
            \end{pmatrix}.
      \]
      Dengan $D_0=1$ dan $D_1=p$, berlaku
      \[
            D_n=pD_{n-1}-qrD_{n-2}\qquad(n\ge2).
      \]
\end{prop}
\begin{proof}
      Ekspansi pada baris pertama memberi $pD_{n-1}$ dari entri $(1,1)$.
      Minor entri $(1,2)$ mempunyai kolom pertama $(r,0,\ldots,0)^T$,
      sehingga determinannya $rD_{n-2}$. Tanda kofaktor posisi $(1,2)$
      negatif, sehingga suku kedua adalah $-qrD_{n-2}$.
\end{proof}

\begin{exam}[Dua akar rekurensi yang dapat berimpit]
      Ambil $p=a+b$, $q=ab$, dan $r=1$ pada $T_n$.
      Tentukan $D_n$, termasuk ketika $a=b$.
      \sumber{\citet[Soal 2.8]{zhang2009}}
\end{exam}
\begin{solution}
      Rekurensinya
      $D_n=(a+b)D_{n-1}-abD_{n-2}$, dengan $D_0=1$ dan $D_1=a+b$.
      Jika $a\ne b$, barisan
      \[
            D_n=\frac{a^{n+1}-b^{n+1}}{a-b}
            =\sum_{k=0}^n a^{n-k}b^k
      \]
      memenuhi rekurensi dan kedua nilai awal, sehingga merupakan jawabannya.
      Bentuk jumlah menunjukkan langsung bahwa untuk $a=b$,
      \[
            D_n=(n+1)a^n\qquad(n\ge1),
      \]
      dan $D_0=1$. Kasus akar berimpit ini tidak memerlukan pembagian oleh $a-b$.
\end{solution}

\begin{exam}[Kasus Fibonacci]
      Untuk $p=q=1$ dan $r=-1$, rekurensinya $D_n=D_{n-1}+D_{n-2}$,
      dengan $D_0=D_1=1$. Dengan konvensi $F_0=0$, $F_1=1$, diperoleh
      $D_n=F_{n+1}$. Khususnya determinan berukuran $5$ bernilai $F_6=8$.
      \sumber{\citet[Soal 2.2]{zhang2009}; rekurensi ditulis untuk ukuran umum}
\end{exam}

\begin{exam}[Determinan matriks pendamping]
      Untuk $a_0,\ldots,a_{n-1}\in\F$, misalkan
      \[
            C=\begin{pmatrix}
                  0      & 1      & 0      & \cdots   & 0        \\
                  0      & 0      & 1      & \ddots   & \vdots   \\
                  \vdots & \vdots & \ddots & \ddots   & 0        \\
                  0      & 0      & \cdots & 0        & 1        \\
                  -a_0   & -a_1   & \cdots & -a_{n-2} & -a_{n-1}
            \end{pmatrix}.
      \]
      Buktikan $\det(tI_n-C)=t^n+a_{n-1}t^{n-1}+\cdots+a_1t+a_0$.
      \sumber{\citet[Soal 2.10]{zhang2009}, dengan konvensi $\det(tI-C)$}
\end{exam}
\begin{solution}
      Tuliskan determinan yang dicari sebagai $P_n(t)$.
      Ekspansi kolom pertama pada $tI-C$ memberi dua kontribusi:
      $t$ kali determinan berukuran $n-1$ dengan koefisien
      $a_1,\ldots,a_{n-1}$, dan kontribusi $a_0$ dari baris terakhir.
      Pada kontribusi kedua, minor segitiga mempunyai diagonal $-1$,
      sehingga tandanya $(-1)^{n+1}(-1)^{n-1}=1$.
      Dengan demikian,
      \[
            P_n(t)=tP_{n-1}(t;a_1,\ldots,a_{n-1})+a_0.
      \]
      Induksi, dimulai dari $P_1(t)=t+a_0$, memberi polinom yang diminta.
      Perhitungan ini murni merupakan ekspansi determinan.
\end{solution}

\section{Determinan Matriks Blok}
\label{det:blok}

Ukuran blok harus sesuai, tetapi blok tidak dapat diperlakukan sebagai
skalar komutatif. Misalnya, bentuk $\det(AD-BC)$ untuk matriks
$\left(\begin{smallmatrix}A&B\\C&D\end{smallmatrix}\right)$ memerlukan
hipotesis tambahan dan urutan faktor yang tepat.

\begin{prop}[Segitiga blok]
      Jika $A$ dan $D$ persegi, maka
      \[
            \det\begin{pmatrix}A&B\\0&D\end{pmatrix}
            =\det\begin{pmatrix}A&0\\C&D\end{pmatrix}
            =\det A\det D.
      \]
\end{prop}
\begin{proof}
      Pada bentuk pertama, setiap suku Leibniz tak nol harus memasangkan
      baris-baris blok bawah dengan kolom-kolom blok kanan.
      Akibatnya, baris blok atas dipasangkan dengan kolom blok kiri.
      Permutasi terurai menjadi dua permutasi dalam masing-masing blok;
      penjumlahannya menghasilkan $\det A\det D$. Bentuk kedua mengikuti
      transpose. Tidak diperlukan invertibilitas blok diagonal.
\end{proof}

\begin{thm}[Komplemen Schur]
      \label{det:schur}
      Untuk $M=\begin{pmatrix}A&B\\C&D\end{pmatrix}$ dengan $A,D$ persegi,
      \[
            \det M=\det A\det(D-CA^{-1}B)\quad\text{jika }A\text{ invertibel}.
      \]
      Jika $D$ invertibel, rumus alternatifnya adalah
      $\det M=\det D\det(A-BD^{-1}C)$.
      \sumber{rumus pertama pada \citet[Soal 2.35(a)]{zhang2009}}
\end{thm}
\begin{proof}
      Kalikan dari kiri dengan matriks segitiga blok berdeterminan $1$:
      \[
            \begin{pmatrix}I&0\\-CA^{-1}&I\end{pmatrix}
            \begin{pmatrix}A&B\\C&D\end{pmatrix}
            =\begin{pmatrix}A&B\\0&D-CA^{-1}B\end{pmatrix}.
      \]
      Gunakan determinan hasil kali dan segitiga blok.
      Untuk rumus kedua, gunakan pengali
      $\left(\begin{smallmatrix}I&-BD^{-1}\\0&I\end{smallmatrix}\right)$.
\end{proof}

\begin{thm}[Kondisi komutativitas yang cukup]
      \label{det:blokkomutatif}
      Jika $A,B,C,D\in\Mat_n(\F)$ dan $AC=CA$, maka
      \[
            \det\begin{pmatrix}A&B\\C&D\end{pmatrix}=\det(AD-CB).
      \]
      \sumber{\citet[Soal 2.35(b)]{zhang2009}}
\end{thm}
\begin{proof}
      Jika $A$ invertibel, rumus Schur memberi
      \begin{align*}
            \det M
             & =\det A\det(D-CA^{-1}B)          \\
             & =\det(AD-ACA^{-1}B)=\det(AD-CB).
      \end{align*}
      Kesamaan terakhir memakai $AC=CA$.
      Untuk $A$ singular, ganti $A$ dengan $A+tI$.
      Matriks ini tetap komutatif dengan $C$ dan invertibel untuk semua
      $t$ kecuali sejumlah hingga nilai, sebab $\det(A+tI)$ adalah
      polinom monik berderajat $n$.
      Kedua determinan menjadi polinom dalam $t$ yang sama pada tak hingga
      banyak nilai, sehingga identik. Substitusi $t=0$ menyelesaikan kasus singular.
\end{proof}

Urutan $CB$ pada kesimpulan harus dipertahankan.
Latihan~\ref{det:lat-blok-syarat} menunjukkan bahwa menukarnya menjadi
$BC$, atau menghapus syarat $AC=CA$, dapat mengubah hasil.

\begin{prop}[Identitas Sylvester untuk determinan]
      Untuk $A\in\Mat_{m\times n}(\F)$ dan $B\in\Mat_{n\times m}(\F)$,
      \[
            \det(I_m+AB)=\det(I_n+BA).
      \]
\end{prop}
\begin{proof}
      Hitung determinan
      $\left(\begin{smallmatrix}I_m&A\\-B&I_n\end{smallmatrix}\right)$
      dengan komplemen Schur terhadap $I_m$, lalu terhadap $I_n$.
      Kedua hasilnya adalah ruas kanan dan ruas kiri identitas tersebut.
\end{proof}

Identitas ini mengurangi ukuran masalah jika salah satu ukuran jauh
lebih kecil. Contohnya, jika $U,V$ berukuran $n\times2$, maka
\[
      \det(I_n+UV^T)=\det(I_2+V^TU).
\]
Dengan demikian, perubahan rank paling tinggi dua dapat dihitung
melalui satu determinan $2\times2$, berapa pun besarnya $n$.

\begin{exam}[Blok dengan pola simetris]
      Untuk sembarang $A,B\in\Mat_n(\F)$, buktikan
      \[
            \det\begin{pmatrix}A&B\\B&A\end{pmatrix}
            =\det(A+B)\det(A-B).
      \]
      \sumber{bagian determinan pada \citet[Soal 2.40]{zhang2009}}
\end{exam}
\begin{solution}
      Ambil $T=\left(\begin{smallmatrix}I&I\\I&-I\end{smallmatrix}\right)$,
      sehingga $T^{-1}=T/2$. Perkalian blok memberi
      \[
            T^{-1}\begin{pmatrix}A&B\\B&A\end{pmatrix}T
            =\begin{pmatrix}A+B&0\\0&A-B\end{pmatrix}.
      \]
      Determinan tidak berubah oleh keserupaan, lalu rumus blok diagonal
      memberikan hasilnya. Syarat $AB=BA$ tidak diperlukan.
\end{solution}

\begin{exam}[Determinan blok real yang tidak negatif]
      Jika $A,B\in\Mat_n(\R)$, buktikan
      \[
            \det\begin{pmatrix}A&B\\-B&A\end{pmatrix}\ge0.
      \]
      \sumber{bagian real dari \citet[Soal 2.39]{zhang2009}}
\end{exam}
\begin{solution}
      Lakukan perhitungan di $\C$ dengan
      $T=\left(\begin{smallmatrix}I&I\\iI&-iI\end{smallmatrix}\right)$.
      Matriks ini invertibel dan
      \[
            T^{-1}\begin{pmatrix}A&B\\-B&A\end{pmatrix}T
            =\begin{pmatrix}A+iB&0\\0&A-iB\end{pmatrix}.
      \]
      Karena $A,B$ real, kedua determinan blok diagonal saling konjugat.
      Oleh karena itu,
      \[
            \det\begin{pmatrix}A&B\\-B&A\end{pmatrix}
            =\det(A+iB)\,\overline{\det(A+iB)}
            =|\det(A+iB)|^2\ge0.
      \]
      Jika $A,B$ boleh kompleks, faktorisasi menjadi $\det(A+iB)\det(A-iB)$
      masih benar, tetapi kesimpulan tidak negatif tidak berlaku umum.
      Sebagai contoh, pada ukuran satu, $A=(i)$ dan $B=(0)$ memberi determinan $-1$.
\end{solution}

\clearpage
\section{Latihan dari Zhang}
\label{det:latihan}

Soal-soal berikut dipilih dari Bab 2 \citet{zhang2009} dan dituliskan
kembali dalam bahasa Indonesia. Nomor soal sumber dicantumkan agar
dapat dicocokkan dengan buku. Jawaban perlu disertai identitas atau
struktur yang digunakan, bukan hanya nilai akhir determinan.

\begin{latihan}\label{det:lat-tigadeterminan}
      Tentukan ketiga determinan berikut, dengan $x\in\F$:
      \[
            \begin{vmatrix}1&2&3\\8&9&4\\7&6&5\end{vmatrix},\qquad
            \begin{vmatrix}1+x&2+x&3+x\\8+x&9+x&4+x\\7+x&6+x&5+x\end{vmatrix},
            \qquad
            \begin{vmatrix}x&x^2&x^3\\x^8&x^9&x^4\\x^7&x^6&x^5\end{vmatrix}.
      \]
      \sumber{\citet[Soal 2.1, hlm.~26; PDF hlm.~45]{zhang2009}}
\end{latihan}

\begin{latihan}\label{det:lat-divisibilitas}
      Jika seluruh entri $A\in\Mat_n(\Z)$ bernilai $1$ atau $-1$,
      buktikan bahwa $2^{n-1}$ membagi $\det A$.
      \sumber{\citet[Soal 2.12, hlm.~29; PDF hlm.~48]{zhang2009}}
\end{latihan}

\begin{latihan}\label{det:lat-paritas}
      Misalkan $A$ berukuran $n\times n$.
      \begin{enumerate}[label=(\alph*)]
            \item Jika $A^T=-A$ dan $n$ ganjil, buktikan $\det A=0$.
            \item Jika $A\in\Mat_n(\R)$ memenuhi $A^2+I_n=0$, buktikan $n$ genap.
            \item Apakah kesimpulan (b) tetap berlaku untuk matriks kompleks?
      \end{enumerate}
      \sumber{\citet[Soal 2.14, hlm.~29; PDF hlm.~48]{zhang2009}}
\end{latihan}

\begin{latihan}\label{det:lat-ortogonal}
      Misalkan $A\in\Mat_n(\C)$ memenuhi $AA^T=I_n$ dan $\det A<0$.
      Tentukan $\det(A+I_n)$ tanpa menggunakan nilai eigen.
      \sumber{\citet[Soal 2.15, hlm.~29; PDF hlm.~48]{zhang2009}}
\end{latihan}

\begin{latihan}\label{det:lat-blok-syarat}
      Untuk $M=\left(\begin{smallmatrix}A&B\\C&D\end{smallmatrix}\right)$
      dengan blok real $2\times2$, berikan contoh tandingan untuk:
      \begin{enumerate}[label=(\alph*)]
            \item mengganti $\det(AD-CB)$ oleh $\det(AD-BC)$ meskipun $AC=CA$;
            \item menggunakan $\det M=\det(AD-CB)$ tanpa syarat $AC=CA$.
      \end{enumerate}
      \sumber{\citet[Soal 2.35(c),(d), hlm.~33; PDF hlm.~52]{zhang2009}}
\end{latihan}

\begin{latihan}\label{det:lat-dependensi}
      Misalkan $A\ne0$ dan $B_1,B_2,B_3,B_4\in\Mat_2(\R)$, dengan
      $A\in\Mat_2(\R)$. Jika
      \[
            \det(A+B_i)=\det A+\det B_i,\qquad i=1,2,3,4,
      \]
      buktikan bahwa $B_1,B_2,B_3,B_4$ bergantung linear atas $\R$.
      \sumber{\citet[Soal 2.36(b), hlm.~33; PDF hlm.~52]{zhang2009}}
\end{latihan}

\begin{latihan}\label{det:lat-permutasiblok}
      Tentukan kedua determinan berikut dalam bentuk terfaktor:
      \[
            \begin{vmatrix}
                  0 & 0 & a_1 & b_1 \\0&0&a_2&b_2\\a_3&b_3&0&0\\a_4&b_4&0&0
            \end{vmatrix},\qquad
            \begin{vmatrix}
                  a_1 & 0 & 0 & b_1 \\0&a_2&b_2&0\\0&b_3&a_3&0\\b_4&0&0&a_4
            \end{vmatrix}.
      \]
      \sumber{\citet[Soal 2.4, hlm.~26; PDF hlm.~45]{zhang2009}}
\end{latihan}

\ifpembahasan
      \clearpage
      \section{Pembahasan Latihan}
      \label{det:pembahasan}

      \paragraph{Latihan~\ref{det:lat-tigadeterminan}.}
      Untuk determinan kedua, kurangkan baris pertama dari baris kedua
      dan ketiga:
      \[
            \det\begin{pmatrix}
                  1+x & 2+x & 3+x \\7&7&1\\6&4&2
            \end{pmatrix}
            =10(1+x)-8(2+x)-14(3+x)=-12(x+4).
      \]
      Kasus $x=0$ langsung memberi determinan pertama, yaitu $-48$.
      Untuk determinan ketiga, ekspansi enam suku memberi
      \begin{align*}
            D(x) & =x^{15}+x^{13}+x^{17}-x^{19}-x^{15}-x^{11} \\
                 & =-x^{11}(1-x^2)(1-x^6).
      \end{align*}
      Seluruh langkah merupakan identitas polinom, sehingga hasilnya juga
      berlaku untuk $x=0$ tanpa pemeriksaan pembagian tersendiri.

      \paragraph{Latihan~\ref{det:lat-divisibilitas}.}
      Untuk setiap $i\ge2$, lakukan $R_i\leftarrow R_i-R_1$.
      Setiap entri pada $n-1$ baris tersebut menjadi $0$, $2$, atau $-2$.
      Keluarkan faktor $2$ dari masing-masing baris. Maka
      $\det A=2^{n-1}\det B$ untuk suatu matriks bilangan bulat $B$.
      Karena $\det B\in\Z$, keterbagian yang diminta terbukti.

      \paragraph{Latihan~\ref{det:lat-paritas}.}
      (a) Transpose dan perkalian skalar memberi
      \[
            \det A=\det(A^T)=\det(-A)=(-1)^n\det A.
      \]
      Jika $n$ ganjil, $2\det A=0$, sehingga determinannya nol.

      (b) Persamaan $A^2=-I_n$ memberi $(\det A)^2=(-1)^n$.
      Karena $A$ real, ruas kiri tidak negatif, sehingga $n$ tidak mungkin
      ganjil. (c) Untuk matriks kompleks kesimpulan itu gagal: matriks
      $1\times1$ dengan entri $i$ sudah memenuhi $A^2+I_1=0$.

      \paragraph{Latihan~\ref{det:lat-ortogonal}.}
      Dari $AA^T=I$ diperoleh $(\det A)^2=1$, sehingga $\det A=\pm1$,
      meskipun entri $A$ boleh kompleks. Syarat negatif memilih $\det A=-1$.
      Selain itu, $A^{-1}=A^T$, maka
      \[
            \det(A+I)=\det A\det(I+A^{-1})
            =-\det(I+A^T)=-\det(I+A).
      \]
      Jadi $\det(A+I)=0$.

      \paragraph{Latihan~\ref{det:lat-blok-syarat}.}
      Untuk (a), ambil
      \[
            A=\begin{pmatrix}1&0\\0&2\end{pmatrix},\quad
            C=\begin{pmatrix}3&0\\0&4\end{pmatrix},\quad
            B=\begin{pmatrix}0&1\\1&0\end{pmatrix},\quad
            D=\begin{pmatrix}1&1\\0&1\end{pmatrix}.
      \]
      Jelas $AC=CA$. Tetapi
      \[
            AD-CB=\begin{pmatrix}1&-2\\-4&2\end{pmatrix},\qquad
            AD-BC=\begin{pmatrix}1&-3\\-3&2\end{pmatrix}.
      \]
      Determinan pertama $-6$, sama dengan $\det M$ menurut
      \autoref{det:blokkomutatif}, sedangkan determinan kedua $-7$.

      Untuk (b), pilih
      \[
            A=\begin{pmatrix}1&1\\0&1\end{pmatrix},\quad
            B=\begin{pmatrix}0&1\\1&0\end{pmatrix},\quad
            C=\begin{pmatrix}1&0\\0&2\end{pmatrix},\quad D=I_2.
      \]
      Maka $AC\ne CA$. Komplemen Schur memberi
      \[
            D-CA^{-1}B=\begin{pmatrix}2&-1\\-2&1\end{pmatrix},
            \qquad \det M=1\cdot0=0.
      \]
      Namun $AD-CB=\left(\begin{smallmatrix}1&0\\-2&1\end{smallmatrix}\right)$
      berdeterminan $1$. Jadi hipotesis komutativitas tidak dapat dihapus.

      \paragraph{Latihan~\ref{det:lat-dependensi}.}
      Tuliskan $A=\left(\begin{smallmatrix}a&b\\c&d\end{smallmatrix}\right)$
      dan $X=\left(\begin{smallmatrix}x&y\\z&w\end{smallmatrix}\right)$.
      Ekspansi langsung memberikan
      \[
            \det(A+X)-\det A-\det X=dx-cy-bz+aw.
      \]
      Jadi keempat vektor koordinat $\operatorname{vec}(B_i)$, dengan urutan
      koordinat $(x,y,z,w)^T$, memenuhi persamaan linear homogen yang sama:
      \[
            (d,-c,-b,a)\operatorname{vec}(B_i)=0.
      \]
      Baris $(d,-c,-b,a)$ tidak nol karena $A\ne0$.
      Jika $K$ adalah matriks $4\times4$ dengan kolom-kolom
      $\operatorname{vec}(B_i)$, maka $(d,-c,-b,a)K=0$ memaksa $K$ singular.
      Kolom-kolom $K$ bergantung linear; hubungan linear yang sama berlaku
      untuk $B_1,B_2,B_3,B_4$.

      \paragraph{Latihan~\ref{det:lat-permutasiblok}.}
      Pada determinan pertama, tukarkan dua kelompok kolom berukuran dua.
      Diperlukan $2\cdot2=4$ pertukaran, sehingga tandanya positif.
      Matriks menjadi blok diagonal dan determinannya
      \[
            (a_1b_2-a_2b_1)(a_3b_4-a_4b_3).
      \]
      Pada determinan kedua, urutkan baris \emph{dan} kolom menjadi
      $(1,4,2,3)$. Kedua faktor tanda saling meniadakan, dan terbentuk dua
      blok diagonal
      $\left(\begin{smallmatrix}a_1&b_1\\b_4&a_4\end{smallmatrix}\right)$
      serta $\left(\begin{smallmatrix}a_2&b_2\\b_3&a_3\end{smallmatrix}\right)$.
      Hasilnya
      \[
            (a_1a_4-b_1b_4)(a_2a_3-b_2b_3).
      \]
\fi
